% Clase 21 — ciclo de histéresis saturado (cualitativo) y curva de primera
% imantación punteada. Remanencia: B con H = 0. Coercitividad: H con B = 0.
% El ciclo es simétrico (curva impar), a diferencia del dibujo del pizarrón.
% Ramas: B = Bs tanh((H -+ Hc)/w), con Bs = 1.3, Hc = 1, w = 0.9.
\begin{tikzpicture}
  \begin{axis}[
    emfig, width=9.5cm, height=7cm,
    axis lines=middle,
    xlabel style={anchor=north west, font=\small},
    ylabel style={anchor=south east, font=\small},
    xmin=-4.2, xmax=4.2, ymin=-1.65, ymax=1.65,
    xtick=\empty, ytick=\empty, grid=none, clip=false,
    xlabel={$H$}, ylabel={$B$},
  ]
    \addplot[figblue, domain=-4:4, samples=120] {1.3*tanh((x+1)/0.9)};
    \addplot[figblue, domain=-4:4, samples=120] {1.3*tanh((x-1)/0.9)};
    \addplot[figgray, dashed, domain=0:3.2, samples=60] {1.3*tanh((x/1.6)^1.5)};
    % flechas de recorrido
    \draw[figvec] (axis cs:2.2,1.28) -- (axis cs:1.2,1.18);
    \draw[figvec] (axis cs:-2.2,-1.28) -- (axis cs:-1.2,-1.18);
    % remanencia y coercitividad
    \fill[figred] (axis cs:0,{1.3*tanh(1/0.9)}) circle (2pt);
    \draw[figred, line width=0.5pt] (axis cs:0,1.04) -- (axis cs:-1.1,1.4);
    \node[font=\footnotesize, text=figred, anchor=east] at (axis cs:-1.1,1.4) {remanencia};
    \fill[figred] (axis cs:-1,0) circle (2pt);
    \draw[figred, line width=0.5pt] (axis cs:-1,0) -- (axis cs:-2.0,-0.45);
    \node[font=\footnotesize, text=figred, anchor=north east] at (axis cs:-2.0,-0.45) {coercitividad};
    \draw[figgray, line width=0.5pt] (axis cs:1.9,0.95) -- (axis cs:2.5,0.45);
    \node[font=\footnotesize, text=figgray, anchor=north west] at (axis cs:2.45,0.45) {primera imantación};
  \end{axis}
\end{tikzpicture}
